% Chapter 5 can be included from a larger thesis with \input or \include.
% Citation keys are maintained in references.bib.

\chapter{Real-Time Music--Motion Matching and Continuous Control}
\label{chap:realtime_matching_control}

\section{Chapter Overview}
\label{sec:ch5_overview}

Chapter~\ref{chap:offline_catalogue} constructed a catalogue in which every
audio reference has a consistent descriptor and every selectable dance has a
preflight-passed Unitree G1 trajectory. This chapter develops the online method
that uses this catalogue. The method receives an audio stream, retrieves
musically related references, ranks their associated motions, accumulates
evidence for a stable selection, and connects the selected trajectory to the
currently playing one. At the same time, a causal beat controller adjusts
motion phase without discontinuously jumping between poses.

The proposed method is retrieval-based rather than generative. It cannot create
an unseen choreography, but it offers three practical advantages for the
present humanoid system. First, the source of every commanded pose is known.
Second, expensive retargeting and preflight are completed before live use.
Third, retrieval, selection, timing, and transition errors can be measured
separately. These properties complement recent work on expressive audio-driven
humanoid locomotion and streaming character generation, which demonstrates the
importance of musical responsiveness but also highlights the difficulty of
causal control \cite{li2026you,ji2026discoforcing}.

The online method is organised around three different decisions:

\begin{enumerate}[label=(\arabic*)]
    \item \emph{what} music and motion are supported by the recent audio.
    \item \emph{when} the evidence is stable enough to authorise a switch.
    \item \emph{how} the current and selected robot trajectories can be joined
          without a pose, velocity, contact, or root discontinuity.
\end{enumerate}



\section{Online Problem Formulation}
\label{sec:ch5_problem}

Let $\mathcal{A}_{\leq t}$ denote audio observed up to current time $t$, and let
$\mathcal{M}=\{M_i\}_{i=1}^{N}$ be the finite catalogue of robot-ready
motions. Each $M_i$ is associated with source music $u_i$, an authored tempo,
an activity profile, key-pose phases, foot-contact states, and a G1 trajectory.
The online system selects motion $M_{i_t}$ and phase $\phi_t$, then produces
the emitted configuration $\mathbf{q}^{\mathrm{out}}_t$ using only
$\mathcal{A}_{\leq t}$ and stored state:

\begin{equation}
    \bigl(i_t,\phi_t,\mathbf{q}^{\mathrm{out}}_t\bigr)
    =\mathcal{G}\!\left(\mathcal{A}_{\leq t},\mathcal{M},\mathcal{H}_{t}\right),
    \label{eq:ch5_causal_mapping}
\end{equation}

where $\mathcal{H}_{t}$ contains the current motion, recent rankings, accepted
beats, played-motion history, prepared candidates, and controller state. The
history term makes explicit that selection is stateful. Applying a memoryless
nearest-neighbour rule independently to every window would make
$i_t$ oscillate when several catalogue items have similar scores.

Throughout the thesis, $\mathcal{F}$ denotes audio feature extraction and
$\mathcal{G}$ denotes stateful online control. The waveform history
$\mathcal{A}_{\leq t}$ is distinct from its descriptor $\mathbf{x}_t$.
The full catalogue $\mathcal{D}$ contains descriptors, trajectories $M_i$ and
metadata $\mathbf{m}_i$; $\mathcal{M}$ is its trajectory collection, and
$\mathcal{C}_t$ is the ranked motion candidate set. Thus $\mathbf{x}_i$ and
$M_i$ denote different parts of an item, rather than alternative names for it.
The index $i_t$ identifies the selected motion, $\phi_t$ is playback phase,
and $\mathbf{q}^{\mathrm{out}}_t$ is the emitted robot configuration after
blending and output constraints. The time subscript refers to the current
update instant; $\Delta t$ is the elapsed time between updates.

The method optimises no single end-to-end objective at runtime. Instead, it
uses a sequence of constrained local objectives:

\begin{equation}
\begin{aligned}
    &\text{audio retrieval:} &&\max s_{\mathrm{track}},\\
    &\text{motion ranking:}  &&\max s_{\mathrm{final}},\\
    &\text{entry search:}    &&\min C_{\mathrm{entry}},\\
    &\text{playback:}        &&\min |e_{\phi}|\quad
       \text{subject to phase-rate and output constraints}.
\end{aligned}
    \label{eq:ch5_local_objectives}
\end{equation}

This factorisation reflects the available evidence. Audio descriptors can
indicate semantic and rhythmic compatibility, but they cannot determine the
instantaneous joint configuration at which a switch is smooth. Conversely,
joint-space entry cost cannot decide whether a dance is appropriate for the
music.

\subsection{Dual-timescale causal processing}
\label{subsec:ch5_dual_timescale}

The same audio stream is processed at two timescales. A six-second rolling
waveform supplies the context required by the catalogue extractor. A retrieval
attempt is made once per second when the previous attempt has completed. In
parallel, a short-block analyser emits activity, onset, spectral, and beat
events used for phase control and subtle pose modulation. The retrieval view
therefore estimates musical identity and compatibility, whereas the event view
estimates local timing.

Both paths are causal. No sample later than the current audio position is used.
This is stricter than merely playing an offline model incrementally, a
distinction emphasised in real-time diffusion control
\cite{ji2026discoforcing}. The longer retrieval window does introduce an
observation delay at startup and after a musical change. This delay is explicit
and can be measured. It is not hidden by using future context.

Feature extraction and catalogue search run in a single background worker. If
that worker is occupied, the next submission opportunity is skipped. Requests
are not accumulated in a queue. Thus the live controller consumes the newest
completed evidence instead of processing an increasing backlog of stale audio
windows. Completed rankings are passed to the bounded candidate-preparation
subsystem described in Section~\ref{subsec:ch5_commitment_readiness}.

\section{Hierarchical Music-to-Motion Retrieval}
\label{sec:ch5_hierarchical_retrieval}

\subsection{Query standardisation and segment similarity}
\label{subsec:ch5_segment_similarity}

For every query window, the online extractor produces the same three vectors
as the offline pipeline: a learned embedding $\mathbf{e}_q$, a rhythm--timbre
vector $\mathbf{r}_q$, and style activations $\mathbf{g}_q$. The embedding is
derived from the same Discogs EffNet model used to construct the catalogue
\cite{alonso2022music}. Model identity and feature dimensions are checked at
startup. A mismatch is treated as an error.

Learned and rhythm vectors are standardised using the catalogue statistics and
then normalised to unit length. For feature family $f\in\{e,r\}$,

\begin{equation}
    \overline{\mathbf{f}}_q=
    \frac{(\mathbf{f}_q-\boldsymbol{\mu}_f)
          /\max(\boldsymbol{\sigma}_f,\epsilon)}
         {\left\|(\mathbf{f}_q-\boldsymbol{\mu}_f)
          /\max(\boldsymbol{\sigma}_f,\epsilon)\right\|_2}.
    \label{eq:ch5_query_standardisation}
\end{equation}

Style activations are unit-normalised without z-standardisation because their
coordinates have a direct label meaning. Cosine similarity is mapped into
$[0,1]$:

\begin{equation}
    s_f(q,n)=\frac{1+\overline{\mathbf{f}}_q^{\mathsf{T}}
        \overline{\mathbf{f}}_n}{2}.
    \label{eq:ch5_mapped_cosine}
\end{equation}

For catalogue segment $n$, the fused score is

\begin{equation}
    s_{\mathrm{seg}}(q,n)
    =0.70s_e(q,n)+0.20s_r(q,n)+0.10s_g(q,n).
    \label{eq:ch5_segment_fusion}
\end{equation}

The learned representation supplies the main semantic neighbourhood, while
explicit rhythm prevents two stylistically similar but rhythmically dissimilar
passages from being treated as identical. Style activations have the smallest
weight because they are useful priors but are too coarse to determine a track
on their own. If style-vector dimensions are unavailable, the remaining
weights are renormalised to $0.70/0.90$ and $0.20/0.90$. Missing evidence is
not interpreted as zero similarity.

\subsection{Track pooling and genre-first restriction}
\label{subsec:ch5_track_pooling}

A track normally contributes several overlapping catalogue segments. Using its
single best segment is sensitive to accidental local agreement, whereas an
average over the entire track penalises longer recordings and music with
legitimate internal variation. Let
$s_{u,(1)}\geq s_{u,(2)}\geq\cdots$ be the ordered segment scores belonging to
track $u$. The track score is the mean of at most the three highest values:

\begin{equation}
    s_{\mathrm{track}}(q,u)
    =\frac{1}{K_u}\sum_{k=1}^{K_u}s_{u,(k)},
    \qquad K_u=\min(3,N_u).
    \label{eq:ch5_track_pooling}
\end{equation}

Tracks are next grouped by their AIST++ genre code. An initial genre score is
the same top-three mean over tracks in that genre. Discogs labels are also
mapped to broad AIST++ families. For example, hip-hop labels support the
middle- and low-hip-hop families, while house and techno labels support the
house family. These mappings are deliberately broad priors rather than a claim
that Discogs and AIST++ define identical taxonomies.

When the largest mapped tag prior is at least $0.08$, the final genre score is

\begin{equation}
    s_{\mathrm{genre}}(h)
    =0.35\,\widetilde{s}_{\mathrm{track\mbox{-}genre}}(h)
     +0.65\,\widetilde{p}_{\mathrm{tag}}(h),
    \label{eq:ch5_genre_prior}
\end{equation}

where both terms are normalised within the current query. The default
style-first policy admits the highest-scoring family. It also admits the second
family when the genre-score margin is no more than $0.025$. All later motion
ranking is restricted to admitted families and to tracks among the retrieved
set. This hierarchy reduces an important failure mode of flat retrieval. A
small acoustic advantage from an unrelated genre cannot be compensated later
by a mechanically convenient motion.

\subsection{Abstention and match confidence}
\label{subsec:ch5_abstention}

The system is allowed to abstain. It computes a negative-style score as the
largest activation among non-music, ambient, drone, field-recording, noise, and
related tags. If this value is at least $0.17$ for two consecutive retrieval
windows, the query is rejected as non-dance or ambient material. Persistence
prevents one uncertain tag activation from cancelling a valid dance selection.
A query is also rejected when the beat-quality value defined in
Equation~\ref{eq:beat_quality} is below $0.35$.

For diagnostics, confidence combines beat quality $q_b$ and the leading genre
margin $\Delta_h$:

\begin{equation}
    c_{\mathrm{match}}
    =\operatorname{clip}_{[0,1]}\!\left(
       0.55q_b+0.45\operatorname{clip}_{[0,1]}
       \left(\frac{\Delta_h}{0.10}\right)
      \right).
    \label{eq:ch5_match_confidence}
\end{equation}

This confidence value explains the strength of the evidence but does not
override the explicit rejection rules. When a result is rejected, no new
motion evidence is accumulated and the current motion remains active. Therefore, the
system fails by holding a known trajectory.

\subsection{Motion-specific compatibility}
\label{subsec:ch5_motion_compatibility}

Music retrieval identifies relevant recordings, but each recording can be
associated with several dances that differ in speed, event density, and
activity. A second ranking stage evaluates the motion itself. Only motions
whose offline preflight passed and whose music and genre survived the preceding
stage are considered.

Tempo compatibility handles half-time and double-time ambiguity. For query BPM
$b_q$ and the source tempo $b_i$ of motion $M_i$, the candidate playback ratios
are

\begin{equation}
    \mathcal{R}_i=
    \left\{\frac{b_q}{b_i},\frac{b_q}{2b_i},
                  \frac{2b_q}{b_i}\right\}.
    \label{eq:ch5_tempo_ratios}
\end{equation}

Only ratios in $[0.55,1.60]$ are retained. The valid ratio closest to one in
log space is selected, and its compatibility is

\begin{equation}
    s_{\mathrm{tempo}}(i)
    =\exp\!\left[-\frac{1}{2}
      \left(\frac{\log r_i}{0.28}\right)^2\right].
    \label{eq:ch5_tempo_score}
\end{equation}

If no ratio is feasible, the motion is excluded. If reliable tempo is absent,
the implementation assigns a deliberately weak score of $0.25$ at authored
speed.

Rhythmic compatibility compares the audio event rate with the motion key-pose
rate. The query and motion rates are

\begin{equation}
    \rho_q=0.70\frac{b_q}{60}+0.30d_o,
    \qquad
    \rho_i=r_i d_{k,i},
    \label{eq:ch5_event_rates}
\end{equation}

where $d_o$ is onset density and $d_{k,i}$ is confidence-weighted key-pose
density. Their scale-symmetric agreement and the motion regularity term are

\begin{align}
    s_{\mathrm{density}}(i)
      &=\exp\!\left[-\frac{1}{2}
        \left(\frac{|\log(\rho_i/\rho_q)|}{0.65}\right)^2\right],\\
    s_{\mathrm{key}}(i)
      &=0.80s_{\mathrm{density}}(i)
        +0.20\exp\!\left[-\min(c_{v,i},2)\right].
    \label{eq:ch5_keypoint_score}
\end{align}

When either rate is effectively zero, $s_{\mathrm{key}}=0.25$. This avoids an
undefined logarithm while preserving the uncertainty penalty.

Motion activity is obtained from the catalogue-relative 90th-percentile
velocity. With catalogue median $\mu_v$ and robust scale $\sigma_v$,

\begin{equation}
    a_i=\frac{1}{1+\exp[-(v_{90,i}-\mu_v)/\sigma_v]},
    \qquad
    s_{\mathrm{activity}}(i)=1-|a_q-a_i|,
    \label{eq:ch5_activity_score}
\end{equation}

where $a_q$ is the audio activity from
Equation~\ref{eq:audio_activity}. The complete motion and final scores are

\begin{align}
    s_{\mathrm{motion}}(i)
      &=0.50s_{\mathrm{tempo}}(i)
        +0.35s_{\mathrm{key}}(i)
        +0.15s_{\mathrm{activity}}(i),\\
    s_{\mathrm{final}}(i)
      &=0.65s_{\mathrm{track}}(u_i)
        +0.35s_{\mathrm{motion}}(i).
    \label{eq:ch5_final_motion_score}
\end{align}

Music evidence remains dominant, but the motion-specific term prevents a
semantically similar track from authorising choreography that would require
extreme time scaling or has incompatible event density. All component values
are retained for evaluation and ablation.

\section{Stable Selection and Candidate Readiness}
\label{sec:ch5_stable_selection}

\subsection{Relevance hysteresis}
\label{subsec:ch5_relevance_hysteresis}

Let $i_t^{*}$ be the leading motion at retrieval update $t$, and let $i_c$ be
the currently playing motion. A relevance switch is armed only when

\begin{equation}
    i_t^{*}\neq i_c,
    \qquad
    s_{\mathrm{final}}(i_t^{*})
       \geq s_{\mathrm{final}}(i_c)+0.08,
    \label{eq:ch5_relevance_margin}
\end{equation}

for three consecutive completed retrievals with the same $i_t^{*}$. If the
leader changes, the current motion becomes best, the score margin disappears,
or the match is rejected, the streak is reset. A successful streak creates a
\emph{pending relevance selection}. It does not immediately alter playback.

This hysteresis is important because overlapping audio windows are strongly
correlated. Three wins should not be interpreted as three independent
statistical observations. They are instead a temporal persistence rule that
filters short score inversions while giving sustained evidence a bounded path
to a switch.

\subsection{Diversity under a relevance constraint}
\label{subsec:ch5_diversity_selection}

A stable musical passage can keep the same motion at rank one for a long time.
To avoid visible repetition, diversity is handled by a second policy. It
considers the latest top five motions and retains only candidates satisfying

\begin{equation}
\begin{aligned}
    s_{\mathrm{final}}(i)&\geq s_{\mathrm{final}}(i^{*})-0.05,\\
    s_{\mathrm{track}}(u_i)&\geq s_{\mathrm{track}}(u_{i^{*}})-0.08.
\end{aligned}
    \label{eq:ch5_diversity_eligibility}
\end{equation}

The candidate must remain eligible for the same three-update streak. After the
current motion has been held for four accepted-beat bars, the selector first
prefers a motion cluster absent from the three-item recent history, then an
unseen motion identifier, and finally the least recently played eligible item.
The score constraints are applied before recency, so variation cannot introduce
an arbitrarily poor match.

Relevance and diversity selections record different reasons. A relevance
selection can be committed at the next valid boundary as soon as it is ready.
A diversity selection is committed only after the maximum hold has actually
been reached. This distinction allows a genuine musical change to take
priority over a cosmetic rotation.

\subsection{Accepted-beat commitment and prepared candidates}
\label{subsec:ch5_commitment_readiness}

The switch clock counts accepted alignment beats rather than raw detector
peaks. Every four accepted beats defines a bar boundary. A pending selection
becomes eligible for commitment at this boundary, ensuring that switching and
phase control share one rhythmic interpretation. The logical state progression
is

\begin{equation}
    \texttt{current}\rightarrow\texttt{evidence}
    \rightarrow\texttt{pending}\rightarrow\texttt{prepared}
    \rightarrow\texttt{transitioning}\rightarrow\texttt{current}.
    \label{eq:ch5_selection_states}
\end{equation}

Whenever a ranking arrives, the pending motion and other high-ranked motions
are submitted to a bounded preparation pool. Preparation includes canonical
trajectory loading, grounding in the playback model, joint-name adaptation,
and construction of the entry features used in
Section~\ref{sec:ch5_entry_optimisation}. As a bar boundary approaches, ready
candidates are scored against an estimate of the source phase expected at that
boundary.

The pending identifier remains the preferred choice whenever it is fully
prepared. Entry cost chooses the best phase within that motion. It does not
silently replace a musically preferred motion with a less relevant one. If the
preferred motion is not ready, a different ranked and prepared candidate may be
used as an explicit fallback. If no evidence-backed candidate is ready, the
system continues the current trajectory. File I/O and grounding are never
performed synchronously merely to meet a boundary.

\section{Causal Beat and Phase Control}
\label{sec:ch5_beat_phase_control}

\subsection{Activity gate and beat candidates}
\label{subsec:ch5_beat_candidates}

Microphone startup includes a short calibration from which the 90th percentile
of RMS and onset change estimates the local noise floor. Subsequent audio is
active only when its RMS exceeds both a fixed floor and a multiple of the
calibrated floor. A short release interval prevents the activity state from
chattering near the threshold. When audio is inactive, beat history is reset
and motion amplitude returns gradually towards zero.

Beat candidates are produced by predominant local pulse (PLP) estimation,
following Grosche and M\"uller \cite{grosche2010extracting}.
The implementation uses the librosa library \cite{mcfee2015librosa}, applying
\texttt{librosa.beat.plp} to a rolling history of already acquired audio.
PLP combines locally estimated periodicities
into a pulse curve; its peaks supply beat candidates. Block-level onset and
loudness histories supply local event contrast. Analysis can revise estimates
inside its available window, but the controller acts only when results are
delivered and cannot revise emitted poses. This online PLP path is distinct
from the dynamic-programming beat tracker used for catalogue descriptors
\cite{ellis2007beat}. In the present method,
each candidate has pulse confidence $c_p$ and recent-event contrast $c_x$. Its
selection score is

\begin{equation}
    c_b=(1-w_x)c_p+w_xc_x,
    \qquad w_x=0.5.
    \label{eq:ch5_beat_selection_score}
\end{equation}

Candidates below the hard confidence threshold $c_p<0.35$ are rejected before
they can influence phase or the bar counter.

\subsection{Motion-aware beat acceptance}
\label{subsec:ch5_beat_acceptance}

Dense music may contain subdivisions faster than the current dance can travel
between authored key poses. The controller therefore derives an admissible
time interval from the next key-pose gap $T_k$ and the playback-speed limits:

\begin{equation}
    T_{\min}=\eta\frac{T_k}{r_{\max}},
    \qquad
    T_{\max}=\frac{T_k}{r_{\min}},
    \label{eq:ch5_beat_time_window}
\end{equation}

where $\eta$ is the key-point interval ratio and is one in the current matcher
configuration. A beat earlier than $T_{\min}$ after the previous accepted beat
is rejected. When the detector itself proposes periods shorter than this
feasible interval, recent candidate scores define an adaptive quantile
threshold. Weak subdivisions are rejected early in the admissible window. The
threshold falls as elapsed time approaches $T_{\max}$. At the deadline, an
eligible beat is accepted so that the controller cannot wait indefinitely for
an ideal accent.

This policy treats beat detection and beat use as different problems. The
analyser may report all plausible pulses, but only events compatible with the
motion's next authored interval become control targets. Beat-synchronised dance
methods likewise demonstrate that temporal correspondence must be represented
explicitly rather than inferred from visual plausibility alone
\cite{huang2024beat}.

\subsection{Target rate and gradual phase correction}
\label{subsec:ch5_phase_correction}

For accepted beat index $k$, the controller selects the corresponding authored
key-pose phase $\phi_k^{*}$. The wrapped phase error is

\begin{equation}
    e_{\phi}=\operatorname{wrap}_{[-1/2,1/2)}
             \left(\phi_k^{*}-\phi\right).
    \label{eq:ch5_phase_error}
\end{equation}

The first accepted beat uses the full error as a correction request. Later
events use gain $0.25+0.35c_b$, so a high-quality beat creates a stronger but
still gradual correction. In adaptive mode, if $\Delta\phi_k$ is the authored
phase gap and $\Delta t_k$ is the last accepted-beat interval, the target phase
rate is

\begin{equation}
    \omega^{*}=
    \operatorname{clip}_{[r_{\min}\omega_a,r_{\max}\omega_a]}
    \left(\frac{\Delta\phi_k}{\Delta t_k}\right),
    \label{eq:ch5_target_phase_rate}
\end{equation}

where $\omega_a$ is the authored phase rate. The running base rate is
exponentially smoothed towards $\omega^{*}$. Phase correction is not applied by
assigning $\phi\leftarrow\phi_k^{*}$. Instead, the remaining correction $e_r$
is consumed with time constant $\tau_p$:

\begin{equation}
    \Delta\phi_c=e_r\left(1-e^{-\Delta t/\tau_p}\right).
    \label{eq:ch5_desired_phase_correction}
\end{equation}

The correction is clipped so that the effective phase rate remains within the
same speed interval. Its change is also slew-limited. With the resulting
$\omega_{\mathrm{eff}}$, phase evolves continuously as

\begin{equation}
    \phi_{t+1}=\operatorname{wrap}_{[0,1)}
       \left(\phi_t+\omega_{\mathrm{eff}}\Delta t\right).
    \label{eq:ch5_phase_update}
\end{equation}

This mechanism synchronises by temporarily accelerating or decelerating the
authored motion, not by teleporting to a different frame. When beat evidence
times out, the target rate returns to $\omega_a$. When a new motion controller
is created, it inherits the previous tempo, amplitude, beat history, and
effective-rate state, while its phase is initialised at the selected entry.
Thus switching does not restart rhythmic estimation.

\subsection{Bounded pose modulation}
\label{subsec:ch5_pose_modulation}

Timing adaptation preserves the selected trajectory as the primary source of
movement. A secondary modulation layer adds small, zero-centred offsets from
low-, mid-, and high-frequency energy, rhythm density, off-beat ratio, and beat
accent. In the default subtle mode, the authored pose is not multiplied by a
global amplitude. Instead, bounded offsets are applied mainly to the waist,
shoulders, wrists, and knees. For example, high-frequency activity produces a
small alternating wrist and shoulder detail, while low-frequency energy
produces a small knee component.

This layer is intentionally secondary. It improves immediate responsiveness
between slower retrieval updates without changing motion identity or root
trajectory. Because it can create poses not present in the offline artefact,
its result still passes through the output conditioning described in
Section~\ref{sec:ch5_output_conditioning}.

\section{Transition-Entry Optimisation}
\label{sec:ch5_entry_optimisation}

\subsection{Entry candidates and state features}
\label{subsec:ch5_entry_candidates}

Motion matching and motion-graph methods establish the value of comparing
state features at possible transition points
\cite{kovar2023motion,holden2020learned}. The present method uses this principle
only after music has supplied a small candidate set. For each target key pose,
the search includes its central frame and up to three frames on either side.
The neighbouring candidates receive linearly decreasing salience. This local
expansion can find a materially closer robot configuration without turning the
entry search into an unconstrained scan over the whole clip.

At each candidate, the cached state contains named joint positions and
velocities, joint ranges, grounded root position, root tilt after removing yaw,
root linear velocity, root angular speed, and binary left/right foot contacts.
Only joint names shared by source, target, and the output-limits table are used.
The current source phase is mapped to its nearest source frame.

\subsection{Normalised entry cost}
\label{subsec:ch5_entry_cost}

For source frame $s$ and target candidate $j$, joint-pose and velocity costs are

\begin{align}
    C_q(j)&=\sqrt{\frac{1}{D}\sum_{d=1}^{D}
      \left(\frac{q_{j,d}-q_{s,d}}{R_d}\right)^2},\\
    C_v(j)&=\sqrt{\frac{1}{D}\sum_{d=1}^{D}
      \left(\frac{\dot q_{j,d}-\dot q_{s,d}}
      {v_d^{\max}}\right)^2},
    \label{eq:ch5_pose_velocity_costs}
\end{align}

where $R_d$ is the model joint range and $v_d^{\max}$ is the configured speed
limit. Contact cost is the fraction of the two feet whose support states differ:

\begin{equation}
    C_c(j)=\frac{1}{2}\sum_{f\in\{L,R\}}
       \mathbf{1}\!\left[c_{j,f}\neq c_{s,f}\right].
    \label{eq:ch5_contact_cost}
\end{equation}

Root cost averages four separately normalised disagreements:

\begin{equation}
\begin{split}
    C_r(j)=\frac{1}{4}\bigg(&
       \frac{|z_j-z_s|}{0.15}
       +\frac{|\theta_j-\theta_s|}{0.35}
       +\frac{\|\dot{\mathbf{p}}_j-
                    \dot{\mathbf{p}}_s\|_2}{3.0}\\
       &+\frac{|\omega_{r,j}-\omega_{r,s}|}{4\pi}\bigg),
\end{split}
    \label{eq:ch5_root_cost}
\end{equation}

where $\theta$ is non-yaw root tilt. The key-pose term
$C_m(j)=1-s_j$ favours a salient velocity valley. The total entry cost is

\begin{equation}
    C_{\mathrm{entry}}(j)
      =0.45C_q(j)+0.20C_v(j)+0.20C_c(j)
       +0.10C_r(j)+0.05C_m(j).
    \label{eq:ch5_entry_total}
\end{equation}

Pose mismatch receives the largest weight because it is the most immediately
visible discontinuity. Velocity and contact distinguish frames that have a
similar pose but incompatible movement or support. Root state and musical
salience refine the decision without overriding a severe whole-body mismatch.

\subsection{Feasibility and transition duration}
\label{subsec:ch5_entry_feasibility}

An entry with low instantaneous cost can still be unsuitable if it occurs near
the end of a one-shot target. The conservative remaining playback time is

\begin{equation}
    T_{\mathrm{remain}}(j)
      =\frac{(1-\phi_j)T_i}{r_{\max}}.
    \label{eq:ch5_remaining_time}
\end{equation}

It must cover both the transition and at least one remaining bar under the
default configuration. The unquantised transition lower bound is derived from
joint displacement and the same dynamics limits used by the output layer:

\begin{equation}
    T_{\mathrm{raw}}(j)=\max_d\left(
       \frac{1.5|q_{j,d}-q_{s,d}|}{v_d^{\max}},
       \sqrt{\frac{6|q_{j,d}-q_{s,d}|}{a_d^{\max}}}
    \right).
    \label{eq:ch5_transition_lower_bound}
\end{equation}

After enforcing a minimum, the duration is rounded upward to a half-beat
multiple:

\begin{equation}
    T_{\mathrm{tr}}(j)=
    \operatorname{clip}_{[0.35,1.20]}
    \left(
      \left\lceil\frac{\max(0.35,T_{\mathrm{raw}}(j))}{T_b/2}\right\rceil
      \frac{T_b}{2}
    \right).
    \label{eq:ch5_transition_quantisation}
\end{equation}

The candidate is eligible when
$T_{\mathrm{remain}}\geq T_{\mathrm{tr}}+T_{\mathrm{bar}}$. Candidates whose
raw requirement exceeds the maximum transition duration are avoided whenever
a feasible alternative exists. Eligible candidates are ordered first by total
cost, then by salience cost, remaining duration, and phase. If none leaves the
required duration, the earliest key-pose candidate is used as a controlled
fallback. This fallback is recorded. It is not presented as an optimal entry.

\section{Continuous Transition and Root Motion}
\label{sec:ch5_continuous_transition}

\subsection{Target-root alignment}
\label{subsec:ch5_root_alignment}

Catalogue trajectories use clip-local roots. Before a target is blended, its
entry root is aligned to the current world state in horizontal position and
yaw. Let $(\mathbf{p}_s,R_s)$ be the target clip's entry root and
$(\mathbf{p}_c,R_c)$ the current root. With $Y(R)$ denoting the yaw-only
component,

\begin{equation}
    A=Y(R_c)Y(R_s)^{-1}.
    \label{eq:ch5_yaw_alignment}
\end{equation}

For every target frame $(\mathbf{p},R)$, the aligned horizontal root and
orientation are

\begin{align}
    \mathbf{p}'_{xy}
       &=\mathbf{p}_{c,xy}+
          \left[A(\mathbf{p}-\mathbf{p}_s)\right]_{xy},\\
    R'&=AR.
    \label{eq:ch5_root_alignment_transform}
\end{align}

The grounded vertical coordinate and the target roll and pitch are preserved.
The same rigid transform is retained throughout the transition, so target
displacement is preserved rather than repeatedly forcing its root onto the
source.

\subsection{Smooth pose and orientation blending}
\label{subsec:ch5_transition_blending}

For elapsed transition fraction $u\in[0,1]$, the blend weight is the cubic
smoothstep

\begin{equation}
    h(u)=3u^2-2u^3.
    \label{eq:ch5_smoothstep}
\end{equation}

Joint positions and root translation use linear interpolation with $h(u)$,

\begin{equation}
    \mathbf{q}(u)=(1-h(u))\mathbf{q}_c(u)
                   +h(u)\mathbf{q}_n(u),
    \label{eq:ch5_joint_blend}
\end{equation}

while root orientation uses shortest-path quaternion spherical interpolation.
Both the old and new motion controllers continue to advance during the blend.
Consequently, the transition connects two moving trajectories rather than
freezing their entry frames.

After blending, a persistent root-continuity transform maps the active
clip-local displacement into one world trajectory. On a source change, the
world anchor becomes the last emitted world root and the new clip origin
becomes its current raw root. In continuous mode,

\begin{equation}
    \mathbf{p}_{w,xy}(t)
      =\mathbf{a}_{xy}
       +\left[A_a(\mathbf{p}(t)-\mathbf{p}_0)\right]_{xy},
    \label{eq:ch5_world_root_continuity}
\end{equation}

where $\mathbf{a}$ and $A_a$ are the stored world anchor. This prevents each
new catalogue clip from returning the robot to the dataset origin.

\subsection{One-shot completion and state transfer}
\label{subsec:ch5_atomic_transition}

The catalogue motions are treated as one-shot trajectories. Phase is clamped
before it can sample the implicit last-to-first interpolation used by generic
preview samplers. When a clip approaches its end, a prepared transition can be
started even if no regular bar switch is ready. Its blend progress is coupled
to the remaining source phase so that the old trajectory is fully released
before the terminal pose stalls. If no prepared, ranked candidate exists, the
last pose is held. An optional very small stationary breathing signal can avoid
a perfectly frozen appearance without moving the root.

At blend completion, the target sampler, phase controller, phase tracker, root
references, and motion identifier are committed together. Only after this
transfer does the selection policy reset its held-bar count and record
the new item in recency history. This avoids a mixed frame in which, for
example, the new identifier is paired with the old controller state.

\section{Output Conditioning}
\label{sec:ch5_output_conditioning}

\subsection{Stateful joint dynamics limits}
\label{subsec:ch5_joint_limiter}

Entry optimisation and smoothstep blending reduce discontinuity, but music
modulation and two simultaneously moving clips can still produce an excessive
joint target increment. A stateful limiter therefore operates on the final
joint target $q_d$. For every joint $d$, the emitted velocity and acceleration
must satisfy

\begin{equation}
    |\dot q_d|\leq v_d^{\max},
    \qquad
    |\dot q_d-\dot q_d^{-}|\leq a_d^{\max}\Delta t.
    \label{eq:ch5_output_constraints}
\end{equation}

The limiter computes the velocity needed to approach $q_d$, clips it by speed,
and then clips its change around the previous emitted velocity. When a target
stops, a braking-distance test begins deceleration early enough to avoid
overshoot. If the target can be reached within both limits, it is snapped only
to the feasible target rather than passing beyond it. Per-joint limits may be
loaded from a validated table. Otherwise configured defaults are
used. Limiter activation, requested speed, emitted speed, and emitted
acceleration are retained as diagnostics.

\subsection{Collision policy and state consistency}
\label{subsec:ch5_runtime_collision}

Offline preflight covers each canonical trajectory, but it cannot certify every
pose created by online blending or modulation. The runtime supports collision
checking for modified poses or for every pose. When enabled, collision handling
is applied after the joint limiter because it must inspect the pose that would
otherwise be emitted. If it changes that pose, the limiter is reset to the
actual collision-resolved output. Without this reset, the next update would
compute acceleration from a state that the robot never received.

The current default runtime collision mode is disabled because full geometric
checking has a nontrivial online cost. Experiments that claim online collision
protection must therefore enable and report the selected policy explicitly.
Regardless of this runtime option, all source motions have passed the offline
preflight described in Chapter~\ref{chap:offline_catalogue}. The resulting pose
is finally written by joint name into the G1 MuJoCo model for kinematic state
evaluation \cite{todorov2012mujoco}.

\section{Method Observability and Failure Behaviour}
\label{sec:ch5_observability}

The online trace follows the processing boundaries established above.
Retrieval records acceptance, rejection reason, genre, score margins, component
scores, BPM, and leading tags. Selection and preparation record pending reason,
recent cluster, ready candidates, load failures, and explicit fallback use.
Control and transition records contain accepted-beat time, phase rate, remaining
phase correction, entry-cost components, transition duration, blend progress,
root deltas, limiter intervention, and collision override.

These fields distinguish failures that can look alike in the final animation.
For example, a late accent can be attributed to beat rejection, phase-rate
saturation, transition blending, or the joint limiter. Repetition can be
attributed to weak retrieval evidence, rejection, readiness, or the diversity
clock. The associated fallback is always conservative: the system skips stale
work, excludes failed candidates, and otherwise retains the last valid active
state. It never responds to failure with an arbitrary catalogue item or
synchronous retargeting. The trace supplies the intermediate measurements used
by the experiments in the next chapter.
